\subsubsection{Kosaraju SCC}
% original used `vector<...> adj[n]` which is a VLA (non-standard C++).
%     The code below is kept as reference; for a portable template see PLAN_01 fix 1.5.
\begin{lstlisting}[style=mystyle, language=C++]
// TC: O(N + M)
// usa: encuentra componentes fuertemente conexas (SCC) usando dos DFS: una por orden de salida y otra sobre el grafo transpuesto
void dfs(int node, vector<pair<int,int>> adj[], vector<bool> &vis,
stack<int> &st){
	vis[node]=1;
	for(auto [adjN,adjW] : adj[node]){
		if(!vis[adjN]){
			dfs(adjN , adj , vis , st);
		}
	}
	st.push(node);
}

void dfs2(int node, vector<pair<int,int>> adj[], vector<bool> &vis,
vector<int> &scc, int scc_number){
	vis[node]=1;
	scc[node]=scc_number;
	for(auto [adjN,adjW] : adj[node]){
		if(!vis[adjN]){
			dfs2(adjN , adj , vis , scc , scc_number);
		}
	}
}

void solve(){
	int n,m; cin >> n >> m;
	vector< pair<int,int> > adj[n];
	for(int i=0 ; i<m ; i++){
		int u,v,w; cin >> u >> v >> w;
		u--; v--;
		adj[u].emplace_back(v,w);
	}

	vector<bool> vis(n);
	stack<int> st;
	for(int i=0 ; i<n ; i++){
		if(!vis[i]){
			dfs(i , adj , vis , st);
		}
	}

	vector<pair<int,int>> adjT[n];
	for(int i=0 ; i<n ; i++){
		vis[i]=0;
		for(auto [adjN,adjW] : adj[i]){
			adjT[adjN].push_back({i,adjW});
		}
	}

	vector<int> scc(n);
	int scc_number=0;
	while(!st.empty()){
		int node=st.top();
		st.pop();
		if(!vis[node]){
			scc_number++;
			dfs2(node, adjT, vis, scc, scc_number);
		}
	}

	vector<pair<int,int>> adjCond[scc_number];
	for(int i=0 ; i<n ; i++){
		for(auto [adjN,adjW] : adj[i]){
			if(scc[i] != scc[adjN]){
				adjCond[scc[i]-1].emplace_back(scc[adjN]-1,adjW);
			}
		}
	}
}
\end{lstlisting}
